% !TeX root = ../../Book4.tex
% 纲要标题：### 5.3 一般单连通曲面的统一化
% 纲要位置：几何学计划纲要.md，第 3648 行。
% 正文状态：未撰写；本文件目前是框架。

% 纲要内容开始
% 
% * 非紧单连通曲面的适应穷竭引理
% 
%   利用相对紧曲面区域、边界曲线的填洞及单连通性建立圆盘穷竭；本引理在本章证明，不假定已掌握一般非紧曲面分类。
% 
% * 相容区域的归一化映照
% * 正规族极限及非退化性
% 
%   单列归一化导数、单叶函数畸变控制以及逆映射紧性的引理；分别排除常值极限并验证极限像，不以“由正规族”略过核心存在步骤。
% 
% * 双曲型与抛物型的区分
% * 紧单连通情形和黎曼球面
% * 从万有覆盖得到一般曲面的商描述
% 
% 纲要内容结束

\section{一般单连通曲面的统一化}
\label{g4:sec:0503}

待撰写

\subsection{非紧单连通曲面的适应穷竭引理}
\label{g4:subsec:0503:01}
待撰写

\subsection{相容区域的归一化映照}
\label{g4:subsec:0503:02}
待撰写

\subsection{正规族极限及非退化性}
\label{g4:subsec:0503:03}
待撰写

\subsection{双曲型与抛物型的区分}
\label{g4:subsec:0503:04}
待撰写

\subsection{紧单连通情形和黎曼球面}
\label{g4:subsec:0503:05}
待撰写

\subsection{从万有覆盖得到一般曲面的商描述}
\label{g4:subsec:0503:06}
待撰写
% 纲要条目排为编号小节标题，并进入目录与书签；说明段仍保留于纲要注释中。
